\documentclass[a4paper,12pt]{article}
\usepackage[italian]{babel}
\usepackage[utf8]{inputenc}
\usepackage[T1]{fontenc}
\usepackage{geometry}
\usepackage{setspace}
\usepackage{listings}
\usepackage{xcolor}

\geometry{margin=2.5cm}
\singlespacing

\title{Verifica di Intelligenza Artificiale - Classe 4AIQ }
\author{Prof. Fedeli Massimo}
\date{}

\definecolor{codegray}{rgb}{0.95,0.95,0.95}

\lstset{
	backgroundcolor=\color{codegray},
	basicstyle=\ttfamily\small,
	breaklines=false,
	frame=single,
	showstringspaces=false,
	tabsize=4,
	language=Python
}

\begin{document}
	
	\maketitle
	
	\textbf{Cognome e Nome:} \underline{\hspace{8cm}} \hfill
	\textbf{Data:} \underline{\hspace{3cm}}
	
	\vspace{0.5cm}
	
	\textbf{Durata:} 2 ore
	
	\vspace{0.5cm}
	
	\section*{Contesto}
	
	Un'amministrazione comunale vuole stimare il \textbf{tempo medio di percorrenza} di un tratto di strada urbana (in minuti) in funzione di alcune variabili osservabili.
	
	In particolare si considerano i seguenti fattori:
	
	\begin{itemize}
		\item numero di \textbf{veicoli} presenti sulla strada durante l'intervallo osservato;
		\item numero di \textbf{semafori} presenti nel tratto stradale;
		\item \textbf{velocità media consentita} sulla strada (km/h).
	\end{itemize}
	
	Si vuole costruire un modello di \textbf{regressione lineare multipla} utilizzando Python e la libreria \texttt{scikit-learn}.
	
	Scrivere il codice in Python e rispondere alle domande di teoria:
	
	\begin{itemize}
		\item Esercizio 1
		\subitem generare dati simulati. 2 pt
		\subitem addestrare un modello di regressione. 2 pt
		\subitem utilizzare il modello per fare previsioni. 1 pt
		\subitem visualizzare una sezione del piano di regressione in un grafico 3D. 1  pt
		\item Esercizio 2
		\subitem Rispondere alle domande di teoria. 4 pt
	\end{itemize}
	
	\section*{Esercizio 1}
	
	Completa il seguente programma Python scrivendo il codice mancante a partire dai commenti.
	
	\begin{lstlisting}
		# Importare le librerie necessarie:
		# numpy
		# matplotlib.pyplot
		# LinearRegression da sklearn.linear_model
		# Axes3D da mpl_toolkits.mplot3d
		
		
		# -------------------------
		# 1) DATI (simulati)
		# -------------------------
		
		# fissare il seed del generatore casuale
		
		# definire il numero di osservazioni (n)
		
		# generare i valori casuali della variabile numero di veicoli
		# (distribuzione uniforme tra 100 e 900)
		
		# generare i valori casuali della variabile numero di semafori
		# (distribuzione uniforme tra 2 e 15)
		
		# generare i valori casuali della variabile velocita media consentita
		# (distribuzione uniforme tra 30 e 70 km/h)
		
		# generare il rumore casuale usando una distribuzione normale
		
		# definire la formula del tempo di percorrenza usando
		# veicoli, semafori, velocita e rumore
		
		
		# creare la matrice delle variabili indipendenti X
		
		# definire la variabile target y
		
		
		# -------------------------
		# 2) REGRESSIONE
		# -------------------------
		
		# creare il modello di regressione lineare
		
		# addestrare il modello con fit()
		
		# estrarre i coefficienti del modello
		
		# estrarre l'intercetta
		
		# calcolare il coefficiente di determinazione R^2
		
		# stampare l'equazione del modello e i parametri stimati
		
		
		# -------------------------
		# 3) PREVISIONI
		# -------------------------
		
		# definire due nuovi casi di input
		# (veicoli, semafori, velocita)
		
		# utilizzare il modello per effettuare le previsioni
		
		# stampare le previsioni in modo leggibile
		
		
		# -------------------------
		# 4) VISUALIZZAZIONE 3D
		# -------------------------
		
		# fissare un valore di velocita
		
		# creare la figura e gli assi 3D
		
		# selezionare i punti con velocita vicina al valore fissato
		
		# creare uno scatter plot dei punti selezionati
		
		# creare una griglia di valori per veicoli e semafori
		
		# generare le matrici con meshgrid
		
		# calcolare i valori del piano di regressione
		
		# disegnare la superficie del piano
		
		# impostare le etichette degli assi
		
		# impostare il titolo del grafico
		
		# mostrare il grafico
	\end{lstlisting}
	
	\newpage
	
	\section*{Esercizio 2 - Domande teoriche}
	
	\begin{enumerate}
		
		\item Spiega il significato dell'istruzione che imposta il \textit{seed} del generatore casuale.
		
		\item Che cosa rappresentano i coefficienti del modello di regressione?
		
		\item Che cosa rappresenta l'intercetta del modello?
		
		\item Indica almeno due possibili variabili aggiuntive che potrebbero influenzare il tempo di percorrenza in una strada urbana.
		
		\item Descrivi il funzionamento di un modello di regressione lineare.
		
	\end{enumerate}
	
\end{document}